\documentclass[11pt]{amsart}

\usepackage[T1]{fontenc}
\usepackage{lmodern}
\usepackage{microtype}
\usepackage{amsmath,amssymb,amsthm}
\usepackage[hidelinks]{hyperref}

\hypersetup{
  pdftitle={A Simultaneous Exponential Basis for Two Subsets of a Finite Abelian Group}
}

\newtheorem{theorem}{Theorem}
\newtheorem{proposition}[theorem]{Proposition}
\theoremstyle{remark}
\newtheorem{remark}[theorem]{Remark}

\newcommand{\C}{\mathbb{C}}
\newcommand{\rank}{\operatorname{rank}}
\newcommand{\tr}{\operatorname{tr}}

\title[A simultaneous exponential basis]
{A Simultaneous Exponential Basis for Two Subsets of a Finite Abelian Group}

\date{}

\subjclass[2020]{Primary 43A40, 42C15; Secondary 05B35}
\keywords{Finite abelian group, exponential Riesz basis, tight frame,
matroid intersection}

\begin{document}

\begin{abstract}
We answer Question~1.8 of Ferguson, Mayeli, and Sothanaphan
affirmatively: every two nonempty subsets of equal cardinality in a
finite abelian group admit a simultaneous exponential Riesz basis.
More generally, two equal-norm tight frames with a common index set
have a common set of indices, of cardinality the smaller frame
dimension, selecting vectors that are linearly independent in both
frames. The proof is a direct application of Edmonds' matroid
intersection theorem.
\end{abstract}

\maketitle

\section{The result}

Let $G$ be a finite abelian group with character group $\widehat G$.
For $E\subseteq G$ and $\Lambda\subseteq\widehat G$, write
\[
  \mathcal F(E,\Lambda)
  =
  \bigl(\chi(x)\bigr)_{x\in E,\,\chi\in\Lambda}.
\]
Ferguson, Mayeli, and Sothanaphan asked whether every pair of
nonempty equal-cardinality subsets of $G$ has a simultaneous
exponential Riesz basis \cite[Question~1.8]{FMS}:
\begin{quote}
Is it true that for every finite abelian group $G$, and every pair of
nonempty subsets $E_{1}, E_{2}$ of $G$ such that $|E_{1}|=|E_{2}|$,
there exists a collection of exponentials
$\Lambda\subseteq\widehat{G}$ such that the restrictions of the
exponentials to $E_{1}$ form a Riesz basis for $E_{1}$ and the
restrictions to $E_{2}$ form a Riesz basis for $E_{2}$?
\end{quote}
The question was open as of the latest version of \cite{FMS} (arXiv
v6, 19 January 2021), which records that ``the above question for two
subsets is unresolved''; to our knowledge it has not been resolved
since.

\begin{theorem}\label{thm:main}
Let $E_1,E_2\subseteq G$ be nonempty and satisfy
$|E_1|=|E_2|=k$. Then there is a set
$\Lambda\subseteq\widehat G$, with $|\Lambda|=k$, such that both
$\mathcal F(E_1,\Lambda)$ and $\mathcal F(E_2,\Lambda)$ are
nonsingular. Equivalently, the characters in $\Lambda$ restrict to a
Riesz basis of $L^2(E_i)$ for $i=1,2$.
\end{theorem}

For cyclic $G$ the existence of a simultaneous basis for an arbitrary
family of subsets of equal size is due to Kolountzakis; see
\cite[Prop.~7.3]{FMS}, where the observation is credited to him and
proved by taking $\Lambda=\{0,1,\ldots,k-1\}$ under the canonical
identification of $\widehat G$ with $G$ and evaluating a Vandermonde
determinant. Theorem~\ref{thm:main} removes the cyclicity hypothesis,
at the cost of restricting to two subsets.

\section{A common set of frame indices}

Write $[n]=\{1,\ldots,n\}$. For a matrix
$X\in\C^{k\times n}$ and $A\subseteq[n]$, let $X_A$ denote the
submatrix formed by the columns indexed by $A$.

\begin{proposition}\label{prop:frames}
For $i=1,2$, let
\[
  X_i=
  \begin{bmatrix}
    x_{i1}&\cdots&x_{in}
  \end{bmatrix}
  \in\C^{k_i\times n}
\]
satisfy
\[
  X_iX_i^*=\alpha_i I_{k_i},
  \qquad
  \|x_{ij}\|_2^2=\frac{\alpha_i k_i}{n}
  \quad (j\in[n]),
\]
where $\alpha_i>0$. Then there is a set $B\subseteq[n]$, with
$|B|=\min\{k_1,k_2\}$, such that the columns of both $X_{1,B}$ and
$X_{2,B}$ are linearly independent.
\end{proposition}

\begin{proof}
Let $M_i$ be the linear matroid on $[n]$ represented by the columns
of $X_i$, and write
\[
  r_i(A)=\rank X_{i,A}.
\]
Since $X_iX_i^*=\alpha_i I_{k_i}$, the rank of $M_i$ is $k_i$.
For $A\subseteq[n]$, set
\[
  Q_i(A)=X_{i,A}X_{i,A}^*.
\]
Since
\[
  X_iX_i^*-Q_i(A)
  =X_{i,[n]\setminus A}X_{i,[n]\setminus A}^*\succeq0,
\]
we have
\[
  0\preceq Q_i(A)\preceq X_iX_i^*=\alpha_i I_{k_i}.
\]
The matrix $Q_i(A)$ has rank $r_i(A)$, and each of its nonzero
eigenvalues is at most $\alpha_i$. Hence
\[
  \alpha_i r_i(A)
  \geq \tr Q_i(A)
  =\sum_{j\in A}\|x_{ij}\|_2^2
  =\frac{\alpha_i k_i|A|}{n},
\]
so
\begin{equation}\label{eq:density}
  r_i(A)\geq\frac{k_i|A|}{n}.
\end{equation}

Assume without loss of generality that $k_1\leq k_2$. Edmonds'
matroid intersection theorem \cite{Edmonds} gives
\[
  \max\bigl\{|I|: I\text{ is independent in both }M_1\text{ and }M_2\bigr\}
  =
  \min_{A\subseteq[n]}
  \bigl(r_1(A)+r_2([n]\setminus A)\bigr).
\]
Using \eqref{eq:density}, for every $A\subseteq[n]$ we have
\[
  r_1(A)+r_2([n]\setminus A)
  \geq
  \frac{k_1|A|+k_2(n-|A|)}{n}
  \geq k_1.
\]
Thus there is a common independent set of size $k_1$, which is the
largest possible size in $M_1$.
\end{proof}

\begin{remark}\label{rem:density}
The inequalities \eqref{eq:density} are the uniform-density
conditions for the two column matroids: a loopless matroid on $n$
elements, of rank $r$ and with rank function $A\mapsto r(A)$, is
\emph{uniformly dense} if $r(A)\geq r|A|/n$ for every subset $A$ of
its ground set, equivalently if the density $\rho(A)=|A|/r(A)$ never
exceeds the density $n/r$ of the ground set itself. This is
condition~(2) of
\cite[Theorem~1.1]{DevriendtMulas}. The proof above applies verbatim
to show that two uniformly dense matroids on a common ground set have
a common independent set whose size is the smaller of their ranks.

The passage from an equal-norm tight frame to a uniformly dense column
matroid is the complex counterpart of one of the applications of that
theorem in \cite{DevriendtMulas}, where a real representable matroid
is shown to admit a projection representation with constant diagonal
equal to the reciprocal of its density precisely when it is strictly
uniformly dense in the sense used there.
Proposition~\ref{prop:frames} is thus not offered as new in isolation;
the new content is the resolution of \cite[Question~1.8]{FMS}.
\end{remark}

\section{Proof of the simultaneous-basis theorem}

\begin{proof}[Proof of Theorem~\ref{thm:main}]
Set $n=|G|=|\widehat G|$. For $E\subseteq G$, $|E|=k$, let
\[
  F_E=
  \bigl(\chi(x)\bigr)_{x\in E,\,\chi\in\widehat G}
  \in\C^{k\times n}.
\]
Character orthogonality gives
\[
  (F_EF_E^*)_{x,y}
  =
  \sum_{\chi\in\widehat G}\chi(x)\overline{\chi(y)}
  =n\,\delta_{x,y},
\]
and therefore $F_EF_E^*=nI_k$. Every column of $F_E$ has squared
norm $k$, since $|\chi(x)|=1$. Thus $F_E$ satisfies
Proposition~\ref{prop:frames} with $\alpha=n$.

Apply the proposition to $F_{E_1}$ and $F_{E_2}$, with their columns
indexed by $\widehat G$. It yields a set
$\Lambda\subseteq\widehat G$ of cardinality $k$ for which both
$\mathcal F(E_1,\Lambda)$ and $\mathcal F(E_2,\Lambda)$ have
linearly independent columns. These square matrices are therefore
nonsingular. In finite dimension every basis is a Riesz basis.
\end{proof}

\begin{remark}
The theorem is sharp with respect to the number of subsets:
three two-element subsets of $\mathbb Z_2^2$ need not admit a
simultaneous basis \cite[Example~7.2]{FMS}, where this case is
likewise credited to Kolountzakis. The argument above is
qualitative and gives no uniform bound on the condition numbers of
the two Fourier submatrices. It also applies verbatim to any two
row restrictions of equal size in a complex Hadamard matrix normalized
by $HH^*=nI_n$.
\end{remark}

\begin{thebibliography}{9}

\bibitem{DevriendtMulas}
K.~Devriendt and R.~Mulas,
\emph{Uniform density in matroids, matrices and graphs},
Australas. J. Combin. \textbf{94} (2026), no.~1, 145--176.

\bibitem{Edmonds}
J.~Edmonds,
\emph{Matroid intersection},
Ann. Discrete Math. \textbf{4} (1979), 39--49,
\href{https://doi.org/10.1016/S0167-5060(08)70817-3}
{doi:10.1016/S0167-5060(08)70817-3}.

\bibitem{FMS}
S.~Ferguson, A.~Mayeli, and N.~Sothanaphan,
\emph{Riesz bases of exponentials and multi-tiling in finite abelian
groups},
\href{https://arxiv.org/abs/1904.04487}
{arXiv:1904.04487v6} (2021).

\end{thebibliography}

\end{document}